\section{\texorpdfstring{$\Lom$}{L1}: Operational Semantics}
\label{sec:l1}

$\Lzm$ said what a read must return, not how a process could know, without
asking anyone, that a version is complete. That is $\Lom$: a labelled
transition system whose rules are the mechanism the runtime executes.

Each rule is annotated with the source location implementing it; those
annotations are the refinement map from $\Lim$ to $\Lom$.
\Cref{sec:l1-refinement} discharges $\Lom$ against $\Lzm$.

% ===========================================================================
\subsection{Configurations}
\label{sec:l1-config}

\begin{definition}[Local state]\label{def:local-state}
The state of process $\proc{p}$ is a quadruple $(\clk{p}, K_p, \Vers{p}, \Pend{p})$:
\begin{itemize}[leftmargin=1.5em,itemsep=1pt]
  \item $\clk{p} \in \Time$, the \gls{local-clock} (\cref{ts:clock});
  \item $K_p : \Ranks \to \Time$, the \glspl{known-time} (\cref{ts:known}), written
    $\kn{p}{q}$;
  \item $\Vers{p} : \Time \pto \Status \times (\Addrof{p} \to \Bytes)$, the
    partial map of live \glspl{version}, with components $\Vst{p}{\eff}$ and
    $\Vbuf{p}{\eff}$; we set $\Vst{p}{\eff} = \Absent$ when
    $\eff \notin \operatorname{dom} \Vers{p}$;
  \item $\Pend{p}$, a finite multiset of pending requests.
\end{itemize}
\end{definition}

\begin{definition}[Requests]\label{def:requests}
A pending request is one of
\[
\begin{array}{ll}
  \mathsf{lput}(\eff, a, v)            & \text{local put, not yet applied}\\
  \mathsf{lget}(\eff, a, \ell)         & \text{local get, awaiting publication}\\
  \mathsf{iput}(\eff, a, v)            & \text{received put, not yet applied}\\
  \mathsf{iget}(\eff, a, \ell, r, \iota) & \text{received read request from } r\\
  \mathsf{oput}                        & \text{outgoing put, awaiting send completion}\\
  \mathsf{oget}(\eff, \iota)           & \text{outgoing read, awaiting reply}
\end{array}
\]
We write $\opeff{\varrho}$ for the \gls{effective-time} recorded in request
$\varrho$. The kinds mirror \srcrefr{request.h}{33}{40} exactly.
\end{definition}

\begin{definition}[Messages]\label{def:messages}
$\Net$ is a finite multiset of in-flight messages, each of one of the five
forms of \srcrefr{msg.h}{39}{43}:
\[
\begin{array}{ll}
  \mathsf{PUT}(p \to q,\, t,\, \delta,\, a,\, v) &
  \mathsf{GETREQ}(p \to q,\, t,\, \delta,\, a,\, \ell,\, \iota) \\
  \mathsf{GETREP}(q \to p,\, \eff,\, v,\, \iota) &
  \mathsf{HB}(p \to q,\, t) \\
  \mathsf{TOK}(p \to q,\, \ell) &
\end{array}
\]
Every message carries a \emph{stamp}: the value $t$ for $\mathsf{PUT}$,
$\mathsf{GETREQ}$ and $\mathsf{HB}$, the value $\eff$ for $\mathsf{GETREP}$,
and $\ell$ for $\mathsf{TOK}$. We write $\mathit{stamp}(m)$ for it.
\end{definition}

\begin{definition}[Global configuration]\label{def:configuration}
$\Cfg \defeq \bigl\langle (\clk{p}, K_p, \Vers{p}, \Pend{p})_{p \in \Ranks},\
\Net \bigr\rangle$. The initial configuration $\Cfg_0$ has $\clk{p} = 0$,
$\kn{p}{q} = 0$, $\Vers{p} = \emptyset$, $\Pend{p} = \emptyset$ and
$\Net = \emptyset$ for all $p, q$.
\end{definition}

\begin{convention}[Notation in rules]\label{conv:rule-notation}
$\Cfg[\,x \mapsto e\,]$ denotes $\Cfg$ with component $x$ replaced by $e$. We
write $\Pend{p} \uplus \{\varrho\}$ for multiset insertion and
$\Net \uplus \{m\}$ for message injection. Rules are read as: if the premises
hold in $\Cfg$, then $\Cfg \steps \Cfg'$ with $\Cfg'$ as described in the
conclusion.
\end{convention}

% ===========================================================================
\subsection{Transition rules}
\label{sec:l1-rules}

\subsubsection*{Clock and publication of progress}

\begin{mathpar}
\inferrule*[left=\rn{Tick}]
  {t > \clk{p}}
  {\clk{p} \mathrel{{:}{=}} t \\\\
   \Net \mathrel{{:}{=}} \Net \uplus
     \{\, \mathsf{HB}(p \to q,\, t) : q \in \nbr{p} \,\}}
\end{mathpar}

\noindent\srcrefr{ops.c}{37}{53} (broadcast), \srcref{ops.c}{603} and
\srcref{ops.c}{672} (the guard, folded into put and get issue). \rn{Tick}
is not a separate API call: it fires as a side effect of issuing an \gls{operation}
whose \gls{origin-time} exceeds the clock.

\subsubsection*{Issuing operations}

\begin{mathpar}
\inferrule*[left=\rn{Put-Local}]
  {\delta \ge 0 \\ q = p}
  {\Pend{p} \mathrel{{:}{=}} \Pend{p} \uplus
     \{\mathsf{lput}(t + \delta,\, a,\, v)\}}

\inferrule*[left=\rn{Put-Remote}]
  {\delta \ge 1 \\ q \in \nbr{p}}
  {\Pend{p} \mathrel{{:}{=}} \Pend{p} \uplus \{\mathsf{oput}\} \\\\
   \Net \mathrel{{:}{=}} \Net \uplus
     \{\mathsf{PUT}(p \to q,\, t,\, \delta,\, a,\, v)\}}
\end{mathpar}

\begin{mathpar}
\inferrule*[left=\rn{Get-Local}]
  {\delta = 0 \\ q = p}
  {\Pend{p} \mathrel{{:}{=}} \Pend{p} \uplus
     \{\mathsf{lget}(t,\, a,\, \ell)\}}

\inferrule*[left=\rn{Get-Remote}]
  {\delta \le -1 \\ q \in \nbr{p} \\ \iota \text{ fresh}}
  {\Pend{p} \mathrel{{:}{=}} \Pend{p} \uplus \{\mathsf{oget}(t+\delta, \iota)\} \\\\
   \Net \mathrel{{:}{=}} \Net \uplus
     \{\mathsf{GETREQ}(p \to q,\, t,\, \delta,\, a,\, \ell,\, \iota)\}}
\end{mathpar}

\noindent\srcrefr{ops.c}{608}{636} and \srcrefr{ops.c}{677}{720}. The
$\delta = 0 \Rightarrow q = p$ side conditions are the \code{-DSTM\_EDOM}
checks at \srcref{ops.c}{598} and \srcref{ops.c}{667}. Note what the rules do
\emph{not} require: neither checks $\delta$ against $\dist{p}{q}$. That is
\cref{res:lookahead,res:staleness}, and it is assumed, not verified
\devmark{dev:lookahead-uncheck}.

\subsubsection*{Message delivery}

Every delivery raises the receiver's known time; this is the only way known
times ever change.

\begin{mathpar}
\inferrule*[left=\rn{Recv}]
  {m \in \Net \\ \mathit{dst}(m) = p \\ \mathit{src}(m) = r}
  {\Net \mathrel{{:}{=}} \Net \setminus \{m\} \\\\
   \kn{p}{r} \mathrel{{:}{=}} \max\bigl(\kn{p}{r},\, \mathit{stamp}(m)\bigr) \\\\
   \rn{Dispatch}(m)}
\end{mathpar}

\noindent\srcrefr{ops.c}{432}{436}. $\rn{Dispatch}$ is:

\begin{mathpar}
\inferrule*[left=\rn{D-Put}]
  {m = \mathsf{PUT}(r \to p,\, t,\, \delta,\, a,\, v)}
  {\Pend{p} \mathrel{{:}{=}} \Pend{p} \uplus \{\mathsf{iput}(t+\delta,\, a,\, v)\}}

\inferrule*[left=\rn{D-Req}]
  {m = \mathsf{GETREQ}(r \to p,\, t,\, \delta,\, a,\, \ell,\, \iota)}
  {\Pend{p} \mathrel{{:}{=}} \Pend{p} \uplus
     \{\mathsf{iget}(t+\delta,\, a,\, \ell,\, r,\, \iota)\}}
\end{mathpar}

\begin{mathpar}
\inferrule*[left=\rn{D-Rep}]
  {m = \mathsf{GETREP}(r \to p,\, \eff,\, v,\, \iota) \\
   \mathsf{oget}(\eff, \iota) \in \Pend{p}}
  {\Pend{p} \mathrel{{:}{=}} \Pend{p} \setminus \{\mathsf{oget}(\eff,\iota)\} \\\\
   \text{the issuing get returns } v}

\inferrule*[left=\rn{D-Beat}]
  {m = \mathsf{HB}(r \to p,\, t)}
  {\text{no further effect}}
\end{mathpar}

\noindent\srcrefr{ops.c}{438}{506}. \rn{D-Beat} has no effect beyond the
known-time update performed by \rn{Recv}: that update is the heartbeat's
entire purpose.

\subsubsection*{Applying writes}

\begin{mathpar}
\inferrule*[left=\rn{Apply}]
  {\varrho \in \Pend{p} \\
   \varrho \in \{\mathsf{lput}(\eff,a,v),\, \mathsf{iput}(\eff,a,v)\} \\
   \Vst{p}{\eff} \in \{\Absent, \Wr\}}
  {\Vst{p}{\eff} \mathrel{{:}{=}} \Wr \\\\
   \Vbuf{p}{\eff} \mathrel{{:}{=}} \Vbuf{p}{\eff}[a \dots a{+}|v|{-}1 \mapsto v] \\\\
   \Pend{p} \mathrel{{:}{=}} \Pend{p} \setminus \{\varrho\}}
\end{mathpar}

\noindent\srcrefr{ops.c}{223}{258} (local) and \srcrefr{ops.c}{282}{315}
(incoming). Creating a version initialises $\Vbuf{p}{\eff}$ to all zeroes
(\cref{conv:no-inherit}), implemented as the \code{calloc} in
\code{dstm\_region\_version\_find\_or\_create}.

\begin{mathpar}
\inferrule*[left=\rn{Apply-Violation}]
  {\varrho \in \Pend{p} \\
   \varrho \in \{\mathsf{lput}(\eff,a,v),\, \mathsf{iput}(\eff,a,v)\} \\
   \Vst{p}{\eff} \in \{\Rd, \Dc\}}
  {\textbf{abort}}
\end{mathpar}

\noindent\srcrefr{ops.c}{244}{253} and \srcrefr{ops.c}{301}{310}.
\rn{Apply-Violation} is not a semantic rule: it is a detector. On a legal
history it is unreachable (\cref{thm:write-closure}), so reaching it proves
that \cref{res:phase}, \cref{res:monotone} or \cref{res:lookahead} was
violated. The implementation deliberately aborts rather than dropping or
applying the write, because both alternatives silently corrupt a value some
reader has already observed.

\subsubsection*{Promotion --- the two rules that matter}

\begin{mathpar}
\inferrule*[left=\rn{Publish}]
  {\Vst{p}{\eff} = \Wr \\
   \nexists\, \varrho \in \Pend{p}.\;
     \varrho \in \{\mathsf{lput},\mathsf{iput}\} \wedge \opeff{\varrho} = \eff \\\\
   \forall q \in \nbr{p}.\;\ \kn{p}{q} + \dist{p}{q} > \eff \\
   \clk{p} \ge \eff}
  {\Vst{p}{\eff} \mathrel{{:}{=}} \Rd}
\end{mathpar}

\noindent\srcrefr{ops.c}{57}{117}: the pending-put scan at
\srcrefr{ops.c}{63}{76}, the neighbour scan at \srcrefr{ops.c}{86}{101}, the
local-clock clause at \srcref{ops.c}{112}.

\begin{mathpar}
\inferrule*[left=\rn{Retire}]
  {\Vst{p}{\eff} = \Rd \\
   \nexists\, \varrho \in \Pend{p}.\;
     \varrho \in \{\mathsf{lget},\mathsf{iget}\} \wedge \opeff{\varrho} = \eff \\\\
   \forall q \in \nbr{p}.\;\ \kn{p}{q} - \dist{p}{q} > \eff \\
   \clk{p} \ge \eff}
  {\Vst{p}{\eff} \mathrel{{:}{=}} \Dc}
\end{mathpar}

\noindent\srcrefr{ops.c}{121}{172}.

\begin{mathpar}
\inferrule*[left=\rn{Evict}]
  {\Vst{p}{\eff} = \Dc}
  {\Vers{p} \mathrel{{:}{=}} \Vers{p} \setminus \{\eff\}}
\end{mathpar}

\noindent\srcrefr{ops.c}{201}{205}.

The two neighbour clauses are the design in one line. A neighbour $q$ believed
to be at $\kn{p}{q}$ can still write as early as $\kn{p}{q} + \dist{p}{q}$ and
still read as far back as $\kn{p}{q} - \dist{p}{q}$. Publication waits for the
first bound to pass $\eff$, retirement for the second: one \gls{distance}
serving as a forward horizon for writes and a backward horizon for reads.

\subsubsection*{Serving reads}

\begin{mathpar}
\inferrule*[left=\rn{Serve-Local}]
  {\mathsf{lget}(\eff, a, \ell) \in \Pend{p} \\ \Vst{p}{\eff} = \Rd}
  {\Pend{p} \mathrel{{:}{=}} \Pend{p} \setminus \{\mathsf{lget}(\eff,a,\ell)\} \\\\
   \text{the issuing get returns } \Vbuf{p}{\eff}[a \dots a{+}\ell{-}1]}
\end{mathpar}

\begin{mathpar}
\inferrule*[left=\rn{Serve-Remote}]
  {\mathsf{iget}(\eff, a, \ell, r, \iota) \in \Pend{p} \\ \Vst{p}{\eff} = \Rd}
  {\Pend{p} \mathrel{{:}{=}} \Pend{p} \setminus \{\mathsf{iget}(\dots)\} \\\\
   \Net \mathrel{{:}{=}} \Net \uplus
     \{\mathsf{GETREP}(p \to r,\, \eff,\, \Vbuf{p}{\eff}[a \dots a{+}\ell{-}1],\, \iota)\}}
\end{mathpar}

\noindent\srcrefr{ops.c}{260}{280} and \srcrefr{ops.c}{317}{357}. Both rules
require status exactly $\Rd$: a $\Wr$ version is incomplete and a $\Dc$ one is
already forfeit. A get whose target version is $\Absent$ simply does not
fire, which is why violating \cref{res:staleness} manifests as a hang rather
than an error \devmark{dev:lookahead-uncheck}.

\subsubsection*{Termination}

The \gls{barrier} is a three-round token ring, summarised as three rules; the
five-state machine implementing them is \srcrefr{barrier.c}{113}{225}.

\begin{mathpar}
\inferrule*[left=\rn{Bar-One}]
  {p \text{ has no } \mathsf{oput}/\mathsf{oget} \text{ pending} \\
   p = 0 \vee \text{first-round token received}}
  {\text{forward first-round token to } p \oplus 1}

\inferrule*[left=\rn{Bar-Two}]
  {\text{first ring closed} \\ \Pend{p} = \emptyset}
  {\mathit{drained}_p \mathrel{{:}{=}} \mathsf{true} \\\\
   \text{forward second-round token to } p \oplus 1}
\end{mathpar}

\begin{mathpar}
\inferrule*[left=\rn{Bar-Exit}]
  {\text{second ring closed}}
  {\text{forward exit token} \\\\ \mathit{drained}_p \mathrel{{:}{=}} \mathsf{false} \\\\
   \text{barrier returns}}
\end{mathpar}

While $\mathit{drained}_p$ holds, \rn{Publish} and \rn{Retire} drop their
neighbour clauses (\srcref{ops.c}{85}, \srcref{ops.c}{149}). This is sound
because the closed first ring certifies no process will issue another
\gls{operation}, making those horizons vacuous; and necessary, since without
it the last \glspl{version} could never publish, no further \glspl{heartbeat}
arriving to raise the known times.

% ===========================================================================
\subsection{Fairness}
\label{sec:l1-fairness}

Safety holds for every execution; liveness needs the usual assumptions for an
asynchronous message-passing system.

\begin{definition}[Admissible execution]\label{def:admissible}
An execution is \emph{admissible} when
\begin{enumerate}[label=(F\arabic*),leftmargin=1.5cm,itemsep=1pt]
  \item \label{f:delivery} every message in $\Net$ is eventually delivered by
    \rn{Recv} (reliable, eventual, no loss or duplication);
  \item \label{f:progress} every rule other than \rn{Tick} that stays
    continuously enabled eventually fires --- in particular a process blocked in
    a get keeps driving the engine (\srcref{ops.c}{724});
  \item \label{f:clock} every process either eventually enters the terminating
    barrier or advances its clock without bound.
\end{enumerate}
\end{definition}

\ref{f:clock} formalises the warning in \cref{sec:exec-clocks} that a silent
process blocks its neighbours. It is a restriction on the \emph{application},
not the runtime.

% ===========================================================================
\subsection{Invariants}
\label{sec:l1-invariants}

\begin{invariant}[Clock soundness]\label{inv:l1-clock}
$\kn{p}{q} \le \clk{q}$ for all $p,q$, at every reachable configuration.
\end{invariant}

\begin{proof}
Initially both sides are $0$. $\kn{p}{q}$ changes only in \rn{Recv}, to
$\mathit{stamp}(m)$ for a message $m$ sent by $q$. Each message form is
stamped at or below $q$'s clock at send time: $\mathsf{PUT}$, $\mathsf{GETREQ}$
and $\mathsf{HB}$ carry the origin time $t$, and \rn{Tick} has already set
$\clk{q} \ge t$; $\mathsf{GETREP}$ carries $\eff = t + \delta \le t \le \clk{q}$
because \glspl{get} have $\delta \le 0$; $\mathsf{TOK}$ carries $\ell$, which by
\cref{res:terminate} bounds every effective time and is passed identically by
all \glspl{rank}. Since $\clk{q}$ is nondecreasing (\cref{ts:clock}), the bound
survives.
\end{proof}

\Cref{prop:conservative} now has teeth: known times are lower bounds, both
guards are monotone in them, so every promotion the runtime makes is one a
perfectly-informed oracle would also make. The runtime is never wrong, only
sometimes late.

% ===========================================================================
\subsection{Safety}
\label{sec:l1-safety}

\begin{theorem}[Write closure]\label{thm:write-closure}
Let $\Hist$ be legal. If $\Vst{p}{\eff} \ge \Rd$ in some reachable
configuration, then no operation issued at or after that point is a \gls{put}
targeting $(p, \eff)$, and $\Pend{p}$ contains no unapplied put for $\eff$.
Consequently \rn{Apply-Violation} is unreachable.
\end{theorem}

\begin{proof}
The second clause is the first premise of \rn{Publish}. For the first,
consider a put $w$ targeting $(p,\eff)$ issued by process $\proc{r}$ after the
promotion.

\emph{Case $r \ne p$.} Then $r \in \nbr{p}$, since \cref{res:lookahead}
requires $r \in \nbr{p}$. The \rn{Publish} premise gives
$\kn{p}{r} + \dist{p}{r} > \eff$, and \cref{inv:l1-clock} gives
$\clk{r} \ge \kn{p}{r}$, hence
\[
  \clk{r} + \dist{r}{p} \;\ge\; \kn{p}{r} + \dist{p}{r} \;>\; \eff,
\]
using (M3) for $\dist{r}{p} = \dist{p}{r}$. By \cref{res:monotone} the origin
time of $w$ is at least $\clk{r}$ as it stands now, and by
\cref{res:lookahead} $\opdelta{w} \ge \dist{r}{p}$. Therefore
$\opeff{w} = \optime{w} + \opdelta{w} \ge \clk{r} + \dist{r}{p} > \eff$, so
$w$ does not target $\eff$. Contradiction.

\emph{Case $r = p$.} The \rn{Publish} premise gives $\clk{p} \ge \eff$. By
\cref{res:monotone}, $\optime{w} \ge \clk{p} \ge \eff$, and
$\opdelta{w} \ge 0$, so $\opeff{w} \ge \eff$ with equality only if
$\optime{w} = \eff$ and $\opdelta{w} = 0$. That residual case is exactly a
local zero-\gls{delta} write at the process's current time, and it is excluded by
\cref{res:phase}: the promotion to $\Rd$ can only have been triggered by a get
on $(p,\eff)$ having been issued, and \cref{res:phase} forbids $p$ from
issuing a put on $(p,\eff)$ after such a get.
\end{proof}

\Cref{thm:write-closure} consumes all four mechanical restrictions, and is
where (M2) and (M3) earn their place. Dropping symmetry breaks the inequality
chain in the first case \devmark{dev:asymmetric}; dropping (M2) makes the
second case's strictness fail for a neighbour at distance zero
\devmark{dev:zero-distance}.

\begin{corollary}[Immutability]\label{thm:immutability}
For a legal history, $\Vbuf{p}{\eff}$ is constant from the moment
$\Vst{p}{\eff}$ reaches $\Rd$ until the version is evicted.
\end{corollary}

\begin{proof}
$\Vbuf{p}{\eff}$ is modified only by \rn{Apply}, whose premise requires
$\Vst{p}{\eff} \in \{\Absent, \Wr\}$. By \cref{thm:write-closure} no such
request exists after publication, and statuses never decrease except through
\rn{Evict}, which removes the version outright.
\end{proof}

\begin{theorem}[Discard safety]\label{thm:discard-safety}
Let $\Hist$ be legal. If $\Vst{p}{\eff} = \Dc$ in some reachable
configuration, then no operation issued at or after that point is a get
targeting $(p,\eff)$.
\end{theorem}

\begin{proof}
Let $g$ be such a get, issued by $r$.

\emph{Case $r \ne p$.} Then $r \in \nbr{p}$ by \cref{res:staleness}. The
\rn{Retire} premise gives $\kn{p}{r} - \dist{p}{r} > \eff$, and
\cref{inv:l1-clock} gives $\clk{r} \ge \kn{p}{r}$, so
$\clk{r} - \dist{r}{p} > \eff$ by (M3). By \cref{res:monotone},
$\optime{g} \ge \clk{r}$; by \cref{res:staleness},
$\opdelta{g} \ge -\dist{r}{p}$. Hence
$\opeff{g} \ge \clk{r} - \dist{r}{p} > \eff$.

\emph{Case $r = p$.} The premise gives $\clk{p} \ge \eff$, and
\cref{res:staleness} forces $\opdelta{g} = 0$ for a local get, so
$\opeff{g} = \optime{g} \ge \clk{p} \ge \eff$. Equality would require
$\optime{g} = \eff = \clk{p}$; but retirement additionally requires no pending
$\mathsf{lget}$ at $\eff$, and by \cref{res:phase} combined with
\cref{thm:write-closure} the version was already published, so any local get
at $\eff$ was issued before retirement and has been served.
\end{proof}

\begin{theorem}[Read soundness: $\Lom \refines \Lzm$]\label{thm:read-soundness}
Let $\Hist$ be a legal history. Every get in $\Hist$ that completes in an
$\Lom$ execution returns exactly the bytes prescribed by
\cref{def:read-result} with respect to $\StoreH$.
\end{theorem}

\begin{proof}
A get on $(q,\eff)$ returns via \rn{Serve-Local} or \rn{Serve-Remote},
both of which read $\Vbuf{q}{\eff}$ while $\Vst{q}{\eff} = \Rd$. Define the
abstraction $\alpha(\Cfg)$ mapping each published or retired version to the
byte function $\Vbuf{q}{\eff}$. We show $\alpha(\Cfg)$ agrees with $\StoreH$
on every published version.

Fix a \gls{cell} $c = \cell{q}{\eff}{a}$ with $\Vst{q}{\eff} = \Rd$. By
\cref{thm:write-closure}, every put in $\Hist$ targeting $c$ has already been
applied by \rn{Apply}. By \cref{res:write-disjoint} there is at most one
such put $w$, so no application overwrote another, and the byte in
$\Vbuf{q}{\eff}$ at offset $a$ is the byte $w$ supplied --- which is
$\StoreH(c)$ by \cref{def:store-h}. If there is no such put, no
\rn{Apply} ever touched offset $a$ of this version, so it retains the
zero from version creation (\cref{conv:no-inherit}), and
$\StoreH(c) = 0$ by the second clause of \cref{def:store-h}.
By \cref{thm:immutability} the value does not subsequently change, so every
get on $c$ returns the same byte.
\end{proof}

% ===========================================================================
\subsection{Liveness}
\label{sec:l1-liveness}

Safety says a completed get is right; liveness says it completes.

\begin{theorem}[Deadlock freedom]\label{thm:deadlock-freedom}
Let $\Hist$ be legal and let the execution be admissible
(\cref{def:admissible}). Then every get completes.
\end{theorem}

\begin{proof}
Suppose not, and consider a get blocked forever on $(q, \eff)$. It is waiting
for \rn{Publish} at $(q,\eff)$, whose premises are of three kinds.

The pending-put clause discharges itself: $\Pend{q}$ is finite, no new put
targeting $\eff$ can be created once $q$'s neighbours have advanced (see
below), and \rn{Apply} is always enabled for a pending put on a
non-published version, so by \ref{f:progress} they all clear.

The local-clock clause $\clk{q} \ge \eff$ discharges by \ref{f:clock}: $q$
either advances past $\eff$ or enters the barrier, and in the barrier
$\mathit{drained}_q$ removes the clause.

The neighbour clause requires, for each $r \in \nbr{q}$, that
$\kn{q}{r} > \eff - \dist{q}{r}$. By \ref{f:clock} and \ref{f:delivery} the
heartbeats of \rn{Tick} raise $\kn{q}{r}$ without bound unless $r$ is
itself permanently blocked. So define the \emph{demand relation}: a process
blocked at effective time $\eff$ on process $\proc{q}$ demands that each
$r \in \nbr{q}$ reach origin time strictly greater than
$\eff - \dist{q}{r}$.

Each step of this relation strictly decreases the demanded time, by at least
$\dist{q}{r} \ge 1$ (M2). A cyclic wait would therefore require a chain
$\eff > \eff_1 > \dots > \eff_k = \eff$, which is impossible. Since $\Ranks$
is finite, the demand chain from any blocked process is a finite acyclic chain
terminating at a process whose demand is vacuous, and that process makes progress
by \ref{f:clock}. Working back along the chain, every blocked process is
eventually released.
\end{proof}

\begin{remark}[Why (M2) is not negotiable]\label{rem:zero-cycle}
The proof consumes (M2) exactly once, essentially. Suppose $\dist{p}{q} = 0$
were permitted with its natural reading, ``visible in the same tick''. Then
$p$'s \rn{Publish} clause at $\eff$ reads $\kn{p}{q} > \eff$ and $q$'s reads
$\kn{q}{p} > \eff$: each must observe the other strictly beyond $\eff$ before
either may publish $\eff$, and neither can advance past $\eff$ until it has
read the other's version at $\eff$. The demand chain closes into a zero-weight
cycle and both block forever --- the smallest instance of the classical
zero-lookahead deadlock \cite{chandy1979} \devmark{dev:zero-distance}.
\end{remark}

\begin{theorem}[Barrier correctness]\label{thm:barrier-correct}
Under \cref{def:admissible} and \cref{res:collective,res:terminate}, the
three-round protocol terminates at every process, and when a process returns from
it, $\Pend{p} = \emptyset$ at every process and $\Net$ contains no
$\mathsf{PUT}$, $\mathsf{GETREQ}$ or $\mathsf{GETREP}$.
\end{theorem}

\begin{proof}[Proof sketch]
First round terminates because each process's outgoing set is finite and
\ref{f:delivery} completes the sends. When the first-round token returns to process $0$,
every process has passed \rn{Bar-One}, so no process will emit a further
$\mathsf{PUT}$ or $\mathsf{GETREQ}$. Second round sets $\mathit{drained}_p$, which
removes the neighbour clauses, so every remaining version publishes and
retires without further external input; every pending $\mathsf{iget}$ is then
served and every $\mathsf{oget}$ receives its reply, so $\Pend{p}$ drains.
When the second-round token returns to process $0$, this has happened everywhere. The exit
round propagates that conclusion.
\end{proof}

\begin{proposition}[Three rounds are necessary]\label{thm:barrier-three-rounds}
No two-round ring protocol establishes the postcondition of
\cref{thm:barrier-correct} at every process.
\end{proposition}

\begin{proof}
Two distinct global predicates must be established in sequence:
$\Phi_1$, ``no process will issue another operation'', and $\Phi_2$, ``no process
holds an incomplete request''. $\Phi_2$ cannot be evaluated locally before
$\Phi_1$ is known globally, because until every process has stopped issuing, a
process with an empty queue may still receive a request. A single ring traversal
establishes one predicate, and establishes it \emph{only at the initiator}:
when the token returns to process $0$, process $0$ knows $\Phi_1$, but process $\proc{i}$ with $i > 0$
knows only that processes $0..i$ had passed when the token crossed it.

Hence traversal one establishes $\Phi_1$ at process $0$, and traversal two both
propagates $\Phi_1$ and collects $\Phi_2$, again terminating the knowledge at
process $0$. A process $\proc{i}$ with $i > 0$ that exited after forwarding the second token would
do so without knowing $\Phi_2$ holds at processes $\proc{j}$ with $j > i$; such a process may still
be driving progress, and tearing down \gls{region} $i$ at that moment orphans any
reply still owed to it. A third traversal is required to announce $\Phi_2$.
\end{proof}

% ===========================================================================
\subsection{Refinement}
\label{sec:l1-refinement}

\begin{corollary}[$\Lom$ implements $\Lzm$]\label{cor:refinement}
For every deterministic program whose executions are legal, every admissible
$\Lom$ execution satisfies \cref{def:correctness}.
\end{corollary}

\begin{proof}
Clause (ii) is \cref{thm:read-soundness}. Clause (i) is
\cref{thm:deadlock-freedom} for the gets and \cref{thm:barrier-correct} for
\gls{termination}.
\end{proof}

The refinement map is the annotation on each rule in \cref{sec:l1-rules}.
Three aspects of $\Lim$ are abstracted away:

\begin{description}[leftmargin=1.1cm,style=nextline,font=\normalfont\bfseries]
\item[The allocation policy.] $\Lom$ fixes only the contract of
  \cref{def:alloc-ops} --- fresh handle, contiguous extent, disjoint from
  every live one, stable until freed --- and \cref{ts:alloc} leaves the choice
  of interval nondeterministic. Any policy resolving that nondeterminism
  refines the specification, so the scheme
  \code{dstm\_region\_malloc} happens to implement is not modelled. It
  determines which addresses a derived structure occupies, never what a
  \gls{cell} contains.

\item[MPI tags, buffers and packing.] The $\iota$ of \rn{Get-Remote} is
  the only piece of the messaging layer $\Lom$ retains, because reply matching
  is semantically visible. Packing, non-blocking send progress and the
  request status DFA of \srcrefr{request.h}{43}{46} refine the atomic
  \rn{Recv} and \rn{Serve} rules by \ref{f:delivery}.

\item[Session communicators.] \code{dstm\_start} duplicates a private
  communicator per \gls{session}, which makes $\Net$ per-session. $\Lom$ models one
  session, so this is exactly the assumption that $\Net$ starts empty.
\end{description}
